\section*{Chapitre \ref{ch:b1:precipitation} --- Précipitation et solubilité}

\begin{solution}{exo:b1:precipitation:1}
\ce{BaSO4(s) <=> Ba^2+ + SO4^2-}, $K_s = [\ce{Ba^2+}][\ce{SO4^2-}] =
10^{-9.97}$. \ce{CaF2(s) <=> Ca^2+ + 2F-}, $K_s = [\ce{Ca^2+}][\ce{F-}]^2 =
10^{-9.84}$. \ce{Ag2CrO4(s) <=> 2Ag+ + CrO4^2-}, $K_s =
[\ce{Ag+}]^2[\ce{CrO4^2-}] = 10^{-11.95}$. \ce{Fe(OH)3(s) <=> Fe^3+ + 3OH-},
$K_s = [\ce{Fe^3+}][\ce{OH-}]^3 = 10^{-38.55}$. Pour $K_s = \num{2.0e-13}$,
$\mathrm{p}K_s = 12.70$.
\end{solution}

\begin{solution}{exo:b1:precipitation:2}
$s = \sqrt{K_s} = 10^{-12.27/2} = \qty{7.3e-7}{mol/L}$, soit
$\num{7.3e-7} \times 187.8 = \qty{1.4e-4}{g/L} = \qty{0.14}{mg/L}$.
\end{solution}

\begin{solution}{exo:b1:precipitation:3}
Après le mélange (volume doublé), $[\ce{Ba^2+}] = \qty{1.0e-3}{mol/L}$ et
$[\ce{SO4^2-}] = \qty{1.0e-4}{mol/L}$~: $Q_s = \num{1.0e-7} > K_s =
\num{1.1e-10}$, le sulfate de baryum précipite. Dans le second cas,
$[\ce{SO4^2-}] = \num{1.0e-7}$ et $Q_s = \num{1.0e-10} < K_s$~: pas de
précipité, mais de justesse.
\end{solution}

\begin{solution}{exo:b1:precipitation:4}
Dans l'eau, $s = 10^{-9.97/2} = \qty{1.0e-5}{mol/L}$. Dans une solution
de sulfate à \qty{0.010}{mol/L} (\cref{prop:b1:precipitation:common-ion}),
$s = 10^{-9.97}/0.010 =
\qty{1.1e-8}{mol/L}$, environ mille fois moins (facteur 970).
\end{solution}

\begin{solution}{exo:b1:precipitation:5}
Dans l'eau, $s = (10^{-9.84}/4)^{1/3} = \qty{3.3e-4}{mol/L}$, soit
\qty{26}{mg/L}. Dans le chlorure de calcium à \qty{0.10}{mol/L},
$[\ce{Ca^2+}] \approx 0.10$ et $[\ce{F-}] = 2s$, donc
$s = \frac12\sqrt{10^{-9.84}/0.10} = \qty{1.9e-5}{mol/L}$.
L'approximation $s \ll 0.10$ est vérifiée à quatre ordres de grandeur
près.
\end{solution}

\begin{solution}{exo:b1:precipitation:6}
Début de précipitation~: $\mathrm{pH} = 14.00 - \frac12(11.25 + \log 0.050) = 14.00 - 4.97 =
9.03$. À \qty{99}{\percent}, $[\ce{Mg^2+}] = \num{5.0e-4}$ et
$\mathrm{pH} = 14.00 - \frac12(11.25 - 3.30) = 10.03$. Sur un axe de pH,
\ce{Mg(OH)2} existe au-dessus de 9.03~; en dessous, le magnésium est
entièrement sous forme \ce{Mg^2+}, à \qty{0.050}{mol/L}.
\end{solution}

\begin{solution}{exo:b1:precipitation:7}
Le bromure d'argent apparaît en premier, pour $[\ce{Ag+}] = 10^{-12.27 + 2} =
10^{-10.27}$, le chlorure d'argent pour $10^{-7.75}$. Quand le chlorure
d'argent apparaît, $[\ce{Br-}] = 10^{-12.27}/10^{-7.75} = 10^{-4.52} =
\qty{3.0e-5}{mol/L}$, soit \qty{0.30}{\percent} du bromure. Avec le
critère usuel de \qty{0.1}{\percent}, la séparation n'est pas
acceptable~: les deux
$\mathrm{p}K_s$ ne diffèrent que de 2.5.
\end{solution}

\begin{solution}{exo:b1:precipitation:8}
Le chromate d'argent se forme quand $[\ce{Ag+}]^2 \times \num{5.0e-3} = 10^{-11.95}$,
$[\ce{Ag+}] = \qty{1.5e-5}{mol/L}$. Alors $[\ce{Cl-}] = 10^{-9.75}/\num{1.5e-5}
= \qty{1.2e-5}{mol/L}$, soit une fraction \num{2.4e-4}
(\qty{0.024}{\percent}) du chlorure~: la couleur rouge apparaît quand le
chlorure est pratiquement tout précipité.
\end{solution}

\begin{solution}{exo:b1:precipitation:9}
$[\ce{CO3^2-}] = [\ce{HCO3-}]K_{a2}/[\ce{H3O+}] = [\ce{HCO3-}]\,10^{\mathrm{pH}
- \mathrm{p}K_{a2}}$, donc $K_s = [\ce{Ca^2+}][\ce{HCO3-}]\,10^{\mathrm{pH} -
\mathrm{p}K_{a2}}$, ce qui est la relation demandée. Avec $[\ce{Ca^2+}] =
[\ce{HCO3-}] = s$~: $s^2 = 10^{-8.30 + 10.33 - 8.30} = 10^{-6.27}$,
$s = \qty{7.3e-4}{mol/L}$ (\qty{73}{mg/L} de \ce{CaCO3}), contre
$10^{-4.15} = \qty{7.1e-5}{mol/L}$ si l'ion carbonate n'était pas une
base.
\end{solution}

\begin{solution}{exo:b1:precipitation:10}
1.~En ajoutant \ce{CO2(g) <=> CO2(aq)}~: $K' = K K_H = 10^{-4.34 - 1.47} =
10^{-5.81}$.
2.~$[\ce{Ca^2+}] = s$ et $[\ce{HCO3-}] = 2s$, donc $4s^3 = K'p$. À
\qty{0.010}{bar}~: $s = (10^{-5.81} \times 0.010/4)^{1/3} =
\qty{1.6e-3}{mol/L}$, \qty{157}{mg/L}. Sous l'air extérieur~:
$s = \qty{5.5e-4}{mol/L}$, \qty{55}{mg/L}. (Vérification~: sous l'air du
sol, $[\ce{CO2(aq)}] = 10^{-3.47}$ et $\mathrm{pH} = 6.37 + \log(3.14\times
10^{-3}/3.39 \times 10^{-4}) = 7.34$~: l'hydrogénocarbonate prédomine
bien.)
3.~L'eau saturée en calcite sous la forte pression de dioxyde de
carbone du sol atteint la voûte de la grotte, où cette pression est plus
faible~: elle perd du dioxyde de carbone, $Q > K'$, et la calcite
précipite là où pend la goutte, anneau après anneau.
% ledger: air:CO2, dfg:CO2_g, dfg:CO2_ao (log KH computed in figdata/bachelor-1/aqueous-constants.py)
\end{solution}

\begin{solution}{exo:b1:precipitation:11}
1.~Début de précipitation à $\mathrm{pH} = 14.00 - \frac12(16.38 - 2.00) = 6.81$.
Redissolution complète quand $[\ce{Zn(OH)4^2-}] = 10^{-2} = 10^{-1.93}[\ce{OH-}]^2$,
$[\ce{OH-}] = 10^{-0.035}$, $\mathrm{pH} = 13.97$. L'hydroxyde existe
entre pH 6.81 et 13.97.
2.~$\mathrm{pH} = 14.00 - 14.45/4 = 10.39$, $s_{\min} = 2 \times
10^{-(16.38 + 1.93)/2} = \qty{1.4e-9}{mol/L}$ (dans ce modèle à deux
espèces).
3.~Droite de pente $-2$ de $(6.81, -2)$ jusqu'au minimum vers
$(10.39, -8.9)$, puis de pente $+2$ jusqu'à $(13.97, -2)$~; la courbe
réelle est arrondie par les autres espèces hydroxo (figure de la
\cref{sec:b1:precipitation:amphoteric}).
\end{solution}

\begin{solution}{exo:b1:precipitation:12}
1.~$\ce{Al(OH)3(s) + OH- <=> Al(OH)4-}$, $K = 10^{-1.23}$~; tout
l'aluminium est dissous lorsque $10^{-3} \leqslant K[\ce{OH-}]$,
c'est-à-dire
$[\ce{OH-}] \geqslant 10^{-1.77}$, $\mathrm{pH} \geqslant 12.23$.
2.~$[\ce{Fe^3+}] = 10^{-38.55}/10^{-3 \times 1.77} = 10^{-33.2}$~: le fer
reste entièrement sous forme de \ce{Fe(OH)3}~; une filtration sépare le
fer (solide) de l'aluminium (filtrat).
3.~L'hydroxyde de magnésium n'est pas amphotère~: en excès d'hydroxyde,
il reste solide avec l'hydroxyde de fer(III), et rien n'est séparé.
\end{solution}

\begin{solution}{pb:b1:precipitation:1}
\textbf{1.} \ce{Fe(OH)3(s) <=> Fe^3+ + 3OH-}, $K_s = [\ce{Fe^3+}][\ce{OH-}]^3$;
\ce{Zn(OH)2(s) <=> Zn^2+ + 2OH-}, $K_s = [\ce{Zn^2+}][\ce{OH-}]^2$;
\ce{Mg(OH)2(s) <=> Mg^2+ + 2OH-}, de même.
\textbf{2.} Les constantes n'ont pas la même forme (puissance 3 ou 2 de
$[\ce{OH-}]$) et les trois ions sont à des concentrations différentes~;
seuls les pH de début de précipitation peuvent être comparés.
\textbf{3.} Au début de la précipitation, $c[\ce{OH-}]^n = K_s$, donc
$n\log[\ce{OH-}] =
-\mathrm{p}K_s - \log c$ et $\mathrm{pH} = \mathrm{p}K_e +
\log[\ce{OH-}] = \mathrm{p}K_e - \frac1n(\mathrm{p}K_s + \log c)$.
\textbf{4.} $14.00 - \frac13(38.55 - 3.00) = 2.15$.
\textbf{5.} $14.00 - \frac12(16.38 - 2.30) = 6.96$.
\textbf{6.} $14.00 - \frac12(11.25 - 1.70) = 9.22$.
\textbf{7.} Le fer(III) (2.15), puis le zinc (6.96), puis le magnésium
(9.22). À pH 2.0, $Q_s = 10^{-3} \times 10^{-36} = 10^{-39} < 10^{-38.55}$
pour le fer, et les autres sont encore plus loin de précipiter~: rien
n'est solide.
\textbf{8.} \qty{1.0e-6}{mol/L}.
\textbf{9.} $14.00 - \frac13(38.55 - 6.00) = 3.15$.
\textbf{10.} $\log[\ce{Fe^3+}] = -38.55 + 3(14.00 - \mathrm{pH}) = 3.45 -
3\,\mathrm{pH}$~: pente $-3$.
\textbf{11.} Jusqu'au début de la précipitation de l'hydroxyde de zinc,
pH 6.96.
\textbf{12.} De 3.15 à 6.96.
\textbf{13.} \ce{Fe^3+ + OH- <=> FeOH^2+}, $\beta_1 =
[\ce{FeOH^2+}]/([\ce{Fe^3+}][\ce{OH-}]) = 10^{11.82}$.
\textbf{14.} Les deux sont égales lorsque $[\ce{OH-}] = 10^{-11.82}$, à
pH 2.18~: \ce{Fe^3+} prédomine en dessous, \ce{FeOH^2+} au-dessus.
\textbf{15.} $10^{11.82 + 3.15 - 14.00} = 10^{0.97} = 9.3$~: à pH 3.15, le
fer dissous vaut 10.3 fois le \ce{Fe^3+} libre, soit
\qty{1.0e-5}{mol/L}~: seuls \qty{99.0}{\percent} sont éliminés. La
première estimation n'était pas sûre.
\textbf{16.} $[\ce{Fe}]_{\mathrm{dissous}} = 10^{3.45 - 3\,\mathrm{pH}}\,
(1 + 10^{\mathrm{pH} - 2.18})$.
\textbf{17.} $\log[\ce{FeOH^2+}] = 11.82 + \log[\ce{Fe^3+}] + \mathrm{pH} -
14.00 = 1.27 - 2\,\mathrm{pH}$~: pente $-2$.
\textbf{18.} $1.27 - 2\,\mathrm{pH} \approx -6.00$ donne 3.64~; en
tenant compte du terme \ce{Fe^3+},
$10^{3.45 - 10.92}(1 + 10^{1.46}) =
\num{1.0e-6}$ à $\mathrm{pH} = 3.64$.
\textbf{19.} $[\ce{Zn^2+}] = \num{5.0e-5}$~: $14.00 - \frac12(16.38 - 4.30) =
7.96$. L'hydroxyde de magnésium ne commence à se former qu'à 9.22~: non.
\textbf{20.} \ce{Zn(OH)2(s) + 2OH- <=> Zn(OH)4^2-}, $K = \beta_4 K_s =
10^{14.45 - 16.38} = 10^{-1.93}$.
\textbf{21.} $[\ce{OH-}]^2 = \num{5.0e-3}/10^{-1.93}$, $\mathrm{pH} =
13.81$. L'opérateur ne doit pas ajouter trop de base~: au-dessus de
pH 9.2, le magnésium précipite avec le zinc, et bien au-delà, le zinc
repasse en solution.
\textbf{22.} $\num{1.0e-3} \times 1000 \times 0.999 \times 106.8 =
\qty{107}{g}$ par mètre cube.
\textbf{23.} $\num{5.0e-3} \times 1000 \times 0.99 \times 99.4 =
\qty{492}{g}$ par mètre cube.
\textbf{24.} Si l'on dépasse d'un coup 6.96, les deux hydroxydes
précipitent dans une même boue~: le zinc est perdu dans le déchet ferreux,
et le déchet est contaminé par le zinc. S'arrêter entre les deux seuils
n'élimine que le fer.
\textbf{25.} \textbf{De pH 3.64 à pH 6.96}~: en dessous de 3.64, plus
de \qty{0.1}{\percent} du fer est encore dissous, surtout sous forme de
\ce{FeOH^2+}~; au-dessus de 6.96, l'hydroxyde de zinc se forme. Négliger
\ce{FeOH^2+} aurait placé la borne inférieure à 3.15.
\end{solution}
